\section{Constructible sheaves on affine Grassmannians and affine flag varieties}
\label{jep104local028}

\subsection{Notation}
\label{jep104local037}

Let $\jepPubBbb{F}$ be an algebraically closed field of characteristic $p>0$. Let $G$ be a connected reductive algebraic group over $\jepPubBbb{F}$, let $B^- \subset G$ be a Borel subgroup, and let $T \subset B^-$ be a maximal torus. Let also $B^+ \subset G$ be the Borel subgroup opposite to $B^-$ (with respect to $T$), and let $U^+$ be its unipotent radical.

We denote by $\boldsymbol{X}:=X^*(T)$ the character lattice of $T$, by $\boldsymbol{X}^\vee := X_*(T)$ its cocharacter lattice, by $\Delta \subset \boldsymbol{X}$ the root system of $(G,T)$, and by $\Delta^\vee \subset \boldsymbol{X}^\vee$ the corresponding coroots. We choose the system of positive roots $\Delta^+ \subset \Delta$ consisting of the $T$-weights in $\mathrm{Lie}(U^+)$, and denote by $\boldsymbol{X}^\vee_+ \subset \boldsymbol{X}^\vee$, resp.~$\boldsymbol{X}^\vee_{++} \subset \boldsymbol{X}^\vee$ the corresponding subset of dominant cocharacters, resp.~of strictly dominant cocharacters. We also denote by $\Delta_{\mathrm{s}} \subset \Delta$ the corresponding subset of simple roots, and set
\[
 \rho = \frac{1}{2} \sum_{\alpha \in \Delta^+} \alpha \quad \in \jepPubBbb{Q} \otimes_{\jepPubBbb{Z}} \boldsymbol{X}.
\]
 
 For any $\alpha \in \Delta_{\mathrm{s}}$ we choose an isomorphism between the additive group $\jepPubBbb{G}_{\mathrm{a}}$ and the root subgroup $U_\alpha$ of $G$ associated with $\alpha$, and denote it $u_\alpha$.
 
 We will assume\footnote{This assumption holds in particular if $G$ is semisimple of adjoint type.} that there exists $\varsigma \in \boldsymbol{X}^\vee$ such that $\langle \varsigma, \alpha \rangle = 1$ for any $\alpha \in \Delta_{\mathrm{s}}$; then we have $\boldsymbol{X}^\vee_{++} = \boldsymbol{X}^\vee_+ + \varsigma$. (Such a cocharacter might not be unique; we fix a choice once and for all.)
 
 Let $W_{\mathrm{f}}$ be the Weyl group of $(G,T)$, and let $W:=W_{\mathrm{f}} \ltimes \boldsymbol{X}^\vee$ be the corresponding (extended) affine Weyl group. For $\lambda \in \boldsymbol{X}^\vee$ we will denote by $t_\lambda$ the associated element of $W$. If $w \in W$ and $w = t_\lambda v$ with $\lambda \in \boldsymbol{X}^\vee$ and $v \in W_{\mathrm{f}}$, we set

 \[
  \ell(w) = \sum_{\substack{\alpha \in \Delta^+ \\ v(\alpha) \in \Delta^+}} |\langle \lambda, \alpha \rangle| + \sum_{\substack{\alpha \in \Delta^+ \\ v(\alpha) \in -\Delta^+}} |1+\langle \lambda, \alpha \rangle|.
 \]
 The restriction of $\ell$ to the semi-direct product $W^{\mathrm{Cox}}$ of $W_{\mathrm{f}}$ with the coroot lattice is the length function for a natural Coxeter group structure, and if we set $\Omega := \{w \in W \mid \ell(w)=0\}$ then multiplication induces a group isomorphism $W^{\mathrm{Cox}} \rtimes \Omega \xrightarrow{\sim} W$.

\subsection{The affine Grassmannian and the affine flag variety}
\label{jep104local030}

For the facts we state here, we refer to~\cite{jep104ref21}.

We set $\jepPubScript{K}:=\jepPubBbb{F}( \hspace{-0.5pt} (z) \hspace{-0.5pt} )$, $\jepPubScript{O} := \jepPubBbb{F}[ \hspace{-0.5pt} [z] \hspace{-0.5pt} ]$, and consider the group ind-scheme $G_{\jepPubScript{K}}$ (denoted $LG$ in~\cite{jep104ref21}) and its group subscheme $G_{\jepPubScript{O}}$ (denoted $L^+ G$ in~\cite{jep104ref21}).

We denote by $I^- \subset G_{\jepPubScript{O}}$ the Iwahori subgroup associated with $B^-$, i.e.~the inverse image of $B^-$ under the morphism $G_{\jepPubScript{O}} \to G$ sending $z$ to $0$. We consider the affine Grassmannian $\mathrm{Gr}$ and the affine flag variety $\mathrm{Fl}$ defined by
\[
\mathrm{Gr} := G_{\jepPubScript{K}}/G_{\jepPubScript{O}}, \quad \mathrm{Fl} := G_{\jepPubScript{K}} / I^-.
\]
We denote by $\pi : \mathrm{Fl} \to \mathrm{Gr}$ the projection morphism.

Any $\lambda \in \boldsymbol{X}^\vee$ defines a point $z^\lambda \in T_{\jepPubScript{K}} \subset G_{\jepPubScript{K}}$, hence a point $L_\lambda := z^{\lambda} G_{\jepPubScript{O}} \in \mathrm{Gr}$. We set
\[
\mathrm{Gr}^\lambda := G_{\jepPubScript{O}} \cdot L_\lambda.
\]
Then $\mathrm{Gr}^\lambda$ only depends on the $W_{\mathrm{f}}$-orbit of $\lambda$. Moreover,
the Bruhat decomposition implies that
\[
\mathrm{Gr} = 
\mathop{\textstyle\bigsqcup}\limits_{\lambda \in \boldsymbol{X}^\vee_+} \mathrm{Gr}^\lambda.
\]
We will denote by $j_\lambda : \mathrm{Gr}^\lambda \to \mathrm{Gr}$ the embedding.

For $\lambda \in \boldsymbol{X}^\vee_+$, we will denote by $P_\lambda \subset G$ the parabolic subgroup of $G$ containing $B^-$ associated with the subset of $\Delta_{\mathrm{s}}$ consisting of those simple roots which are orthogonal to $\lambda$. Then $P_\lambda$ is the stabilizer of $L_\lambda$ in $G$, so that we have a canonical isomorphism $G/P_\lambda \xrightarrow{\sim} G \cdot L_\lambda$. Under this identification, it is known that the map $p_\lambda : \mathrm{Gr}^\lambda \to G/P_\lambda$ sending $x$ to $\lim_{t \to 0} t \cdot x$ (where we consider the $\jepPubBbb{G}_{\mathrm{m}}$-action on $\mathrm{Gr}$ via loop rotation) is a morphism of algebraic varieties, and realizes $\mathrm{Gr}^\lambda$ as an affine bundle over $G/P_\lambda$ (see e.g.~\cite[Lem.~2.3]{jep104ref38}).

It is well known (see e.g.~\cite{jep104ref33} or~\cite[\S 2]{jep104ref38}) that if $\lambda \in \boldsymbol{X}^\vee_+$, then we have
\[
\dim(\mathrm{Gr}^\lambda) = \langle \lambda, 2\rho \rangle = \sum_{\alpha \in \Delta^+} \langle \lambda, \alpha \rangle.
\]
We denote by $\preccurlyeq$ the order on $\boldsymbol{X}^\vee_+$ determined by
\[
\lambda \preccurlyeq \mu \quad \text{ iff $\mu - \lambda$ is a sum of positive coroots.}
\]
Then for $\lambda, \mu \in \boldsymbol{X}^\vee_+$ we have
\[
\mathrm{Gr}^\lambda \subset \overline{\mathrm{Gr}^\mu} \quad \text{iff} \quad \lambda \preccurlyeq \mu.
\]

\subsection{Some categories of sheaves on \texorpdfstring{$\mathrm{Gr}$}{Gr} and \texorpdfstring{$\mathrm{Fl}$}{Fl}}
\label{jep104local031}

We let $\ell$ be a prime number which is different from $p$, and let $\Bbbk$ be either a finite extension of $\jepPubBbb{Q}_\ell$, or the ring of integers in such an extension, or a finite field of characteristic $\ell$. In this paper we will be concerned with the constructible derived categories $D^{\mathrm{b}}_c(\mathrm{Gr},\Bbbk)$ and $D^{\mathrm{b}}_c(\mathrm{Fl},\Bbbk)$ of \'etale $\Bbbk$-sheaves on $\mathrm{Gr}$ and $\mathrm{Fl}$, respectively.
If $K \subset G_{\jepPubScript{O}}$ is a subgroup, we will also denote by $D^{\mathrm{b}}_K(\mathrm{Gr},\Bbbk)$ and $D^{\mathrm{b}}_K(\mathrm{Fl},\Bbbk)$ the (constructible) $K$-equivariant derived category of $\Bbbk$-sheaves on $\mathrm{Gr}$ and $\mathrm{Fl}$, in the sense of Bernstein--Lunts~\cite{jep104ref15}. Each of these categories is endowed with the perverse t-structure, whose heart will be denoted $\mathrm{Perv}(\mathrm{Gr},\Bbbk)$, $\mathrm{Perv}(\mathrm{Fl},\Bbbk)$, $\mathrm{Perv}_K(\mathrm{Gr},\Bbbk)$ and $\mathrm{Perv}_K(\mathrm{Fl},\Bbbk)$ respectively.

\begin{rmk}\label{jep104local026}
\begin{enumerate}
\item
Since $\mathrm{Gr}$ and $\mathrm{Fl}$ are ind-varieties and not varieties, the definition of the categories considered above requires some care; 

see e.g.~\cite[\S 2.2]{jep104ref37} or~\cite[App.]{jep104ref27} for details. 
We will not mention this point in the body of the paper, and simply refer to objects in these categories as complexes of sheaves.
\item
Recall that by \cite{jep104ref36} the category $\mathrm{Rep}(G^\vee_R)$ of algebraic representations of the group scheme $G^\vee_R$ over any Noetherian commutative base ring $R$ of finite global dimension is equivalent to the corresponding category of spherical perverse sheaves on $\mathrm{Gr}_{\jepPubBbb{C}}$ in its analytic topology. 
More restrictive assumptions on the base ring in the present paper come from our need to 
use the Artin--Schreier sheaf (see~Section~\ref{jep104local034}), which is only defined in the context of \'etale sheaves over a variety 
in positive characteristic; this setting yields categories of sheaves with coefficients
as above. Notice however that some constructions involving the Artin--Schreier
sheaf do have an analogue for constructible sheaves in the classical topology (see e.g.~\cite{jep104ref44} for the example of Fourier--Deligne transform). We expect that such a
counterpart of the Whittaker category can also be defined (see \cite[Rem. 10.3.6]{jep104ref08}
for a possible approach); this would allow one to extend our main result
to more general coefficient rings.
\end{enumerate}
\end{rmk}

If $K' \subset K \subset G_{\jepPubScript{O}}$ are subgroups, we will denote by
\[
 \mathsf{For}^{K}_{K'} : D^{\mathrm{b}}_K(\mathrm{Gr},\Bbbk) \to D^{\mathrm{b}}_{K'}(\mathrm{Gr},\Bbbk), \quad \mathsf{For}^{K}_{K'} : D^{\mathrm{b}}_K(\mathrm{Fl},\Bbbk) \to D^{\mathrm{b}}_{K'}(\mathrm{Fl},\Bbbk)
\]
the natural forgetful functors. If $K/K'$ is of finite type, these functors have both a right and a left adjoint, which will be denoted ${}^{*}\mathsf{Ind}_{K'}^K$ and ${}^{!}\mathsf{Ind}_{K'}^K$ respectively. If we write $X$ for $\mathrm{Gr}$ or $\mathrm{Fl}$, these functors can be described explicitly by
\[
 {}^{*}\mathsf{Ind}_{K'}^K(\jepPubEuler{F}) = a_* \bigl( \underline{\Bbbk} \, \widetilde{\boxtimes} \, \jepPubEuler{F} \bigr) \quad \text{and} \quad {}^{!}\mathsf{Ind}_{K'}^K(\jepPubEuler{F}) = a_! \bigl( \underline{\Bbbk} \, \widetilde{\boxtimes} \, \jepPubEuler{F} \bigr)[2(\dim K/ K')],
\]
where $\underline{\Bbbk} \, \widetilde{\boxtimes} \,(-)$ is the functor sending an object $\jepPubEuler{F}$ to the only object in $D^{\mathrm{b}}_K(K \times^{K'} X, \Bbbk)$ whose pullback to $K \times X$ (an object of $D^{\mathrm{b}}_{K \times K'}(K \times X,\Bbbk)$, where $K'$ acts on $K \times X$ via $h \cdot (g,x)=(gh^{-1}, h \cdot x)$) is isomorphic to $\underline{\Bbbk}_K \jepDerivedExternalTensor_\Bbbk \jepPubEuler{F}$.
When $K'=\{1\}$ we will write $\mathsf{For}_K$ for $\mathsf{For}^K_{\{1\}}$.

